% 正封图案：画在 (0,0)–(170,240)（单位 mm，不含出血），被 封面-*.tex 引用
% 颜色：cbg 底色、cink 墨色、cacc 橙色、cline 线条灰（和《拆解宇树》《具身智能求职手册》一套）
% 插图：全书的"七级证据阶梯"（第 2 章）。七块台阶从左下到右上：R 研究、D 演示、T 测试、P 试点、O 订单、
% L 交付、M 规模运营。工业机械臂站在最高的 M 级，四足机器人在 L 级，人形机器人（橙色）还在 P 级往上爬。
\begin{scope}
  % 底色（向外多铺 3 mm 出血）
  \fill[cbg] (-3,-3) rectangle (173,243);
  \begin{scope}
    \clip (-3,-3) rectangle (173,243);
    \draw[cline!35, line width=0.15pt] (-3,-3) grid[step=10] (173,243);
  \end{scope}
  % 顶部小字
  \node[anchor=north west, font=\heiti\fontsize{10.5}{12}\selectfont, text=cink!85] at (16,227) {八十年技术谱系 · 每个产品标出证据阶段};
  \draw[cacc, line width=1.8pt] (16,219.5) -- (42,219.5);
  % 书名
  \node[anchor=north west, inner sep=0, font=\heitiblack\fontsize{92}{96}\selectfont, text=cink] at (13,212) {具身智能};
  % 副书名
  \node[anchor=north west, inner sep=0, font=\heiti\bfseries\fontsize{19}{22}\selectfont, text=cacc] at (16,172) {机器、部件与全产业链};

  % ===== 七级阶梯 =====
  % \step{左下 x}{顶面 y}{宽}{正面色}{顶面色}{字母}{中文}
  \def\sDX{9}\def\sDY{5.5}
  \newcommand{\cstep}[7]{%
    \filldraw[fill=#5, draw=cink, line width=0.6pt, line join=round]
      (#1,#2) -- (#1+#3,#2) -- (#1+#3+\sDX,#2+\sDY) -- (#1+\sDX,#2+\sDY) -- cycle;
    \filldraw[fill=#4, draw=cink, line width=0.6pt, line join=round]
      (#1,24) rectangle (#1+#3,#2);
    \filldraw[fill=cink, draw=cink, line width=0.6pt, line join=round]
      (#1+#3,24) -- (#1+#3+\sDX,24+\sDY) -- (#1+#3+\sDX,#2+\sDY) -- (#1+#3,#2) -- cycle;
    \node[font=\heitiblack\fontsize{15}{16}\selectfont, text=white, anchor=north] at (#1+#3/2,#2-2.2) {#6};
    \node[font=\heiti\fontsize{6.6}{7.5}\selectfont, text=white!90, anchor=north] at (#1+#3/2,#2-9.2) {#7};}
  \foreach \i/\L/\C in {0/R/研究,1/D/演示,2/T/测试,3/P/试点,4/O/订单,5/L/交付,6/M/规模运营} {
    \pgfmathsetmacro{\xx}{16+\i*18}
    \pgfmathsetmacro{\yy}{40+\i*11.5}
    \cstep{\xx}{\yy}{18}{cink!82}{white}{\L}{\C}
  }
  % M 级（x 124–142，顶面 y=109）：工业机械臂
  \begin{scope}[shift={(133,111.5)}]
    \fill[cink] (-5,0) rectangle (5,2.6);
    \draw[cink, line width=3.6pt, line cap=round] (0,2.6) -- (0,10) -- (6,22) -- (17,18);
    \draw[cink, line width=1.8pt, line cap=round] (17,18) -- (19.5,13);
    \draw[cink, line width=1.4pt, line cap=round] (18,11.6) -- (19.5,13) -- (21.4,11.8);
    \foreach \p in {(0,10),(6,22),(17,18)} {\filldraw[fill=cbg, draw=cink, line width=0.8pt] \p circle (1.5);}
  \end{scope}
  % L 级（x 106–124，顶面 y=97.5）：四足
  \begin{scope}[shift={(107.5,100)}, scale=0.33]
    \filldraw[fill=cink, draw=cink, rounded corners=3pt] (0,20) rectangle (46,32);
    \filldraw[fill=cink, draw=cink, rounded corners=2pt] (44,26) rectangle (57,36);
    \fill[cacc] (53,31) circle (1.7);
    \foreach \x in {10,20,34,44} {
      \draw[cink, line width=2.2pt, line cap=round] (\x,21) -- (\x-6,10) -- (\x,0);
      \fill[cink] (\x,0) circle (2.2);
    }
  \end{scope}
  % P 级（x 70–88，顶面 y=74.5）：人形机器人（橙色），一条腿迈向上一级
  \begin{scope}[shift={(80,77)}]
    \fill[cacc] (0,27.5) circle (3.1);
    \fill[cbg] (-1.8,27.6) rectangle (1.8,28.6);
    \draw[cacc, line width=5pt, line cap=round] (0,23) -- (0,13);
    \draw[cacc, line width=2.6pt, line cap=round] (0,22) -- (-4.5,15.5) (0,22) -- (5,17.5) -- (8.5,20);
    \draw[cacc, line width=3pt, line cap=round] (0,13) -- (-2.5,6.5) -- (-2.5,0.4);
    \draw[cacc, line width=3pt, line cap=round] (0,13) -- (6,11) -- (9.5,8.5);
  \end{scope}
  % 一行注释
  \node[anchor=north west, font=\heiti\fontsize{10}{14}\selectfont, text=cink!80, align=left] at (16,140) {机械臂在规模运营，\\四足在交付，\\人形机器人还在试点往上爬。};

  % 作者和截止日
  \draw[cink, line width=0.8pt] (16,16.5) -- (154,16.5);
  \node[anchor=south west, font=\heiti\bfseries\fontsize{15}{18}\selectfont, text=cink] at (16,4.5) {甜菜　著};
  \node[anchor=south east, font=\heiti\fontsize{9}{11}\selectfont, text=cink!80] at (154,5.5) {资料截至 2026 年 11 月};
\end{scope}
